% ============================================================
%  Chapter 1: R1 — Four-Dimensional Lorentzian Causal Spacetime
% ============================================================
\chapter{R1 — 因果时空结构}
\label{ch:r1}

\section{总说：宇宙的舞台}

物理学研究的是"什么在什么上面发生"。R1 定义的是"什么上面"——\textbf{舞台本身}。

R1 的核心物理内容由两条\textbf{等价}的表述给出：
\begin{enumerate}
  \item \textbf{最大传播速度}：真空中任何信号、能量或因果影响的传播速度不超过 $c$。$c$ 在所有惯性系中取相同数值——它是普适常数，不是特定参考系的属性。
  \item \textbf{Lorentz 因果结构}（与 (1) 等价）：时空间隔 $ds^2 = g_{\mu\nu}dx^\mu dx^\nu$ 的符号将事件对分为三类。因为 $ds^2 = c^2dt^2 - |d\mathbf{r}|^2$，类空 $\iff |d\mathbf{r}|/dt > c$——所以"无超光速信号"与"类空事件无因果联系"是同一事实的两种表述。
\end{enumerate}

此外：物理定律在所有坐标系中取相同形式（广义协变性）；度规本身可为动力学量——\textbf{引力即时空弯曲}（$\leftarrow$ 等效原理 + 广义协变性）。

\begin{principlebox}[R1 — 因果时空结构]
  时空是四维 Lorentz 流形 $(\mathcal{M}, g_{\mu\nu})$，3 个空间维度和 1 个时间维度。

  \textbf{等价表述}：
  \begin{itemize}
    \item (1a) 最大传播速度：真空中任何信号传播速度 $\leq c$；$c$ 在所有惯性系中不变。
    \item (1b) Lorentz 因果结构：$ds^2 = g_{\mu\nu}dx^\mu dx^\nu$，其符号决定因果分类。
  \end{itemize}

  \textbf{等价性证明}：$ds^2 = c^2dt^2 - |d\mathbf{r}|^2$。类空 $\iff ds^2 < 0 \iff |d\mathbf{r}|/dt > c \iff$ 超光速 $\iff$ 无因果联系。

  广义协变性 + 等效原理 $\to$ 度规为动力学场 $\to$ 引力 = 时空弯曲。
\end{principlebox}

% ── 变量定义表 ──
\section{变量定义}
\label{sec:r1-vars}

\begin{table}[h]
\centering
\small
\begin{tabular}{p{2cm}p{1.8cm}p{2cm}p{1.5cm}p{3cm}p{4cm}}
\toprule
\textbf{变量} & \textbf{符号} & \textbf{类型} & \textbf{量纲} & \textbf{定义域} & \textbf{物理含义} \\
\midrule
时间坐标 & $t$ & $\mathbb{R}$ & T & $t \in \mathbb{R}$ & 类时坐标，参数化因果序列 \\
空间坐标 & $x^i$ & $\mathbb{R}^3$ & L & $i \in \{1,2,3\}$ & 类空坐标 \\
时空坐标 & $x^\mu$ & $\mathbb{R}^4$ & — & $x^0 \equiv ct$ & 4 维流形上的坐标系 \\
度规张量 & $g_{\mu\nu}$ & $\mathbb{R}^{4\times4}$ 对称 & [1] & $\det(g) < 0$，符号差 $(+,-,-,-)$ & 定义线元与因果结构 \\
线元 & $ds^2$ & $\mathbb{R}$ & L$^2$ & $ds^2 \in \mathbb{R}$ & Lorentz 不变量 \\
真空光速 & $c$ & 常数 & L$\cdot$T$^{-1}$ & $c = 2.998\times10^8$ m/s (exact) & 因果传播速度上限 \\
固有时间 & $\tau$ & $\mathbb{R}_{\ge 0}$ & T & $d\tau^2 \equiv ds^2/c^2 \ge 0$ & 类时世界线的弧长参数 \\
空间维数 & $D$ & $\mathbb{N}^+$ & [1] & $D = 3$（contingent） & 宏观延展空间维度数 \\
惯性参考系 & — & 概念 & — & $g_{\mu\nu} = \eta_{\mu\nu}$，$\Gamma^\lambda_{\mu\nu}=0$ & 无加速、无引力的局域参考系 \\
\bottomrule
\end{tabular}
\caption{R1 变量定义}
\end{table}

% ── 数学表达式 ──
\section{数学表达式}

\subsection{局域惯性系（Minkowski 度规 — 平直时空极限）}
\[
ds^2 = c^2 dt^2 - (dx^1)^2 - (dx^2)^2 - (dx^3)^2
\]

\subsection{一般坐标系（弯曲时空）}
\[
ds^2 = g_{\mu\nu}(x)\, dx^\mu dx^\nu
\]
$g_{\mu\nu}$ 为二阶对称协变张量，10 个独立分量（$D=4 \Rightarrow D(D+1)/2 = 10$）。

\subsection{因果分类}
定义间隔 $\Delta s^2 \equiv g_{\mu\nu}\Delta x^\mu\Delta x^\nu$。按 $\Delta s^2$ 的符号分类：

\begin{table}[h]
\centering
\begin{tabular}{p{3cm}p{2.5cm}p{4.5cm}p{4cm}}
\toprule
\textbf{条件} & \textbf{分类} & \textbf{物理含义} & \textbf{速度表述} \\
\midrule
$\Delta s^2 > 0$ & 类时 (timelike) & 可因果连接；存在惯性系中两事件同地先后发生 & $|d\mathbf{r}|/dt < c$ \\
$\Delta s^2 = 0$ & 类光 (lightlike) & 仅以光速信号可连接 & $|d\mathbf{r}|/dt = c$ \\
$\Delta s^2 < 0$ & 类空 (spacelike) & 无因果联系；任何惯性系中均不同地 & $|d\mathbf{r}|/dt > c$ \\
\bottomrule
\end{tabular}
\caption{因果分类——"最大速度"与"光锥结构"等价}
\end{table}

\subsection{局域性约束}
物理量 $\phi(x)$ 在点 $x$ 的值仅依赖于其在过去光锥 $J^-(x)$ 内的初始数据。

% ── 成立条件 ──
\section{成立条件 (Scope)}

\begin{itemize}
  \item \textbf{有效场论范围}：R1 在能量远低于 Planck 尺度（$E \ll E_P \approx 1.22 \times 10^{19}$ GeV）、距离远大于 Planck 长度（$l \gg l_P$）的范围内精确成立。这是目前所有已确认物理实验的适用域。
  \item \textbf{狭义相对论}：$g_{\mu\nu} \rightarrow \eta_{\mu\nu}$（平直极限，忽略引力）。
  \item \textbf{广义相对论}：$g_{\mu\nu}$ 为动力学场，由 Einstein 场方程（$\leftarrow$ R1 + R2）决定。
\end{itemize}

% ── 失效边界 ──
\section{失效边界 (Boundary)}

\begin{limitbox}[R1 的失效边界]
  \begin{itemize}
    \item \textbf{Planck 尺度} $l_P = \sqrt{\hbar G/c^3} \approx 1.616 \times 10^{-35}$ m：时空本身的量子涨落不可忽略。连续流形假设可能破缺。
    \item \textbf{Planck 时间} $t_P = l_P/c \approx 5.391 \times 10^{-44}$ s：经典因果结构可能不再适用。
    \item \textbf{初始奇点}（大爆炸 $t \rightarrow 0$）：度规退化，广义相对论失效。
  \end{itemize}
  \smallskip
  \textit{注：$c$ 见 R1.2，$\hbar$ 见 R4.2，$G$ 为 Newton 引力常数（由 Einstein 场方程的 Newton 极限确定）。}
\end{limitbox}

% ── 必要性 ──
\section{必要性 (Necessity)}

若移除或弱化 R1 的任一部分：

\begin{table}[h]
\centering
\small
\begin{tabular}{p{4cm}p{7cm}}
\toprule
\textbf{移除/弱化} & \textbf{后果} \\
\midrule
$D=3 \to D \neq 3$ & 无稳定 Kepler 轨道；无 Huygens 原理（D 偶数时）；平方反比律改变 \\
Lorentz 符号差 $\to$ Euclidean & 无因果序；无波传播；量子场论无法定义 \\
光速上限 $c \to$ 无上限 & 因果悖论（封闭类时曲线）；无稳定物质结构 \\
广义协变性 $\to$ 特殊坐标系 & 非惯性系物理定律形式改变；无法描述引力为时空几何 \\
\bottomrule
\end{tabular}
\caption{R1 必要性——移除后的物理后果}
\end{table}

% ── 充分性 ──
\section{充分性 (Sufficiency)}

\textbf{R1 单独提供}：
\begin{itemize}
  \item 平直时空极限 $\to$ 狭义相对论全部运动学（Lorentz 变换、时间膨胀、长度收缩、同时性相对性）
\end{itemize}

\textbf{R1 + R2}（最小作用量）$\to$ 弯曲时空动力学（Einstein 场方程）$\to$ 引力理论

\textbf{R1 + R4}（量子公设）$\to$ 相对论量子场论的因果结构（微观因果性 $\to$ 自旋-统计定理）

% ── 子组件映射 ──
\section{子组件与 frameworks.yaml 映射}

\begin{table}[h]
\centering
\begin{tabular}{p{5cm}p{7cm}}
\toprule
\textbf{子组件} & \textbf{frameworks.yaml ID} \\
\midrule
3+1 时空维度 & \texttt{kernel.spacetime\_dimensionality} \\
光速不变 & \texttt{kernel.light\_speed\_invariance} \\
Lorentz 不变性 + 因果结构 & \texttt{kernel.lorentz\_invariance} \\
等效原理（$m_i=m_g$） & \texttt{kernel.equivalence\_principle} \\
广义协变性 & \texttt{kernel.general\_covariance} \\
\bottomrule
\end{tabular}
\caption{R1 子组件与 frameworks.yaml 的对应关系}
\end{table}

\section{直觉理解：光锥——因果结构的可视化}

\subsection{二维类比：从水波到光锥}

想象你站在平静的湖面上，把一块石头扔进水中央。涟漪以圆形向外扩散，速度固定。当你俯视湖面时，你能画出涟漪到达各处的时刻——形成一个向外扩展的圆锥。

时空中的光以完全类似的方式传播。在三维空间中传播形成了"四维圆锥"。我们只能可视化三维（取消一个空间维度），得到熟悉的\textbf{光锥}：

\begin{center}
\begin{tikzpicture}[scale=1.15]
  % Light cone
  \fill[blue!15] (-1.2,-1.2) -- (0,0) -- (1.2,-1.2);
  \fill[blue!5] (-1.2,1.2) -- (0,0) -- (1.2,1.2);

  % Axes
  \draw[->] (0,-1.5) -- (0,1.5) node[above] {$t$ (时间)};
  \draw[->] (-1.5,0) -- (1.5,0) node[right] {$x$ (空间)};

  % Event
  \fill[red] (0,0) circle (2pt) node[below right] {\small 事件 $O$};

  % Labels
  \node at (0.0,1.0) [rotate=0] {\small 未来光锥};
  \node at (0.0,-1.0) [rotate=0] {\small 过去光锥};
  \node at (1.3,0.7) {\small 类时};
  \node at (0.7,0.3) {\small 类光};
  \node at (1.3,-0.7) {\small 类空 (不可达)};
\end{tikzpicture}
\end{center}

\subsection{光锥的物理意义}

\begin{itemize}
  \item \textbf{过去光锥}：能够\emph{影响}事件 $O$ 的所有事件的集合。任何从外部影响 $O$ 的信号都必须来自过去光锥内部或表面上。
  \item \textbf{未来光锥}：能够\emph{被}事件 $O$ \emph{影响}的所有事件的集合。事件 $O$ 发出的信号只能到达未来光锥内部或表面上。
  \item \textbf{类空区域（Elsewhere）}：与 $O$ 绝对分离的区域。此区域内的事件与 $O$ 之间没有任何可能的因果联系——任何观测者都无法将 $O$ 与此区域内的任何事件列为先后关系。
\end{itemize}

\subsection{为什么光速是上限？}

这不是因为"光有特殊性质"，而是因为时空的几何结构。Lorentz 度规决定了间隔 $ds^2$ 的符号：

\[
ds^2 = c^2 dt^2 - |d\mathbf{r}|^2
\]

\begin{itemize}
  \item 若 $v > c$：$|d\mathbf{r}| > c\,dt \implies ds^2 < 0$——事件变为类空分离，失去了因果联系。这意味着信号的传播速度若超过 $c$，在另一个惯性系看来该信号将从"果"传向"因"——这是因果悖论。
  \item 若 $v = c$：$ds^2 = 0$——这是\emph{仅可能}的无质量粒子（光子、引力波）的世界线。
  \item 若 $v < c$：$ds^2 > 0$——有质量粒子只能在类时世界线上运动。
\end{itemize}

\section{从平直时空到弯曲时空}

R1 的真正魅力在于其一般性。Minkowski 度规只描述了平直（无引力）时空。在更一般的弯曲时空中：

\[
ds^2 = g_{\mu\nu}(x)\,dx^\mu dx^\nu
\]

$g_{\mu\nu}$ 是一个位置的函数，有 10 个独立分量（对称性 $g_{\mu\nu}=g_{\nu\mu}$ 将可能的 16 个分量减至 10 个）。

这就是广义相对论的起点：\textbf{引力不是"力"，而是时空的弯曲}。物质告诉时空如何弯曲（通过 Einstein 场方程——从 R2 推出），弯曲的时空告诉物质如何运动（通过测地线方程）。
